% GENERATED FILE. Edit canonical lessons, then run npm run generate:books.
% canonical-source-hash: fnv1a64:ea30ab4be03fa7b4
% canonical-lessons: SA-C201-yojana, SA-C201-laksyam, SA-C201-upayah, SA-C201-samasya, SA-C201-samadhanam

\chapter{Plan, Goal, Way, Problem, Solution}
\label{ch:sa-plan-goal-way-problem-solution}
\hlchaptermodality{201}

\begin{chapteropening}
\textbf{By the end of this chapter:} \emph{I can say a plan, a goal, a way, a problem and a solution.}

Complete the last lesson of chapter 201: I can say a plan, a goal, a way, a problem and a solution.
\end{chapteropening}

\section[\texorpdfstring{\sk{योजना} (yojanā)}{योजना (yojanā)}]{\sk{योजना} (yojanā) — a plan, a scheme}

\label{lesson:SA-C201-yojana}



\begin{quote}
Before the new one: say the Sanskrit for a discussion, then the Sanskrit for a dispute.
\end{quote}

\subsection*{You'll want to know: \sk{योजना}}
\textbf{\sk{योजना}} — \emph{yojanā} — \textquotedblleft{}a plan, a scheme\textquotedblright{}.

\begin{grammarlens}[title={What you've built}]
One more word for this chapter.
\end{grammarlens}

\subsection*{Your turn}
\begin{itemize}
\raggedright
  \item \emph{Say it:} \emph{yojanā}
  \item \emph{Say it:} \emph{yojanā}, once more
  \item \emph{Say it:} say \emph{yojanā}
  \item \emph{Recall:} say the Sanskrit for a discussion, then the Sanskrit for a dispute, then say \emph{yojanā} again
\end{itemize}

\begin{tcolorbox}[breakable,colback=teal!4,colframe=teal!35!black,title={Before you move on}]
What does \sk{योजना} mean? (\textquotedblleft{}A plan, a scheme\textquotedblright{}.) Say it once more.
\end{tcolorbox}
\section[\texorpdfstring{\sk{लक्ष्यम्} (lakṣyam)}{लक्ष्यम् (lakṣyam)}]{\sk{लक्ष्यम्} (lakṣyam) — a goal, a target}

\label{lesson:SA-C201-laksyam}



\begin{quote}
Before the new one: say the Sanskrit for a dispute, then the Sanskrit for a plan.
\end{quote}

\subsection*{You'll want to know: \sk{लक्ष्यम्}}
\textbf{\sk{लक्ष्यम्}} — \emph{lakṣyam} — \textquotedblleft{}a goal, a target\textquotedblright{}.

\begin{grammarlens}[title={What you've built}]
One more word for this chapter.
\end{grammarlens}

\subsection*{Your turn}
\begin{itemize}
\raggedright
  \item \emph{Say it:} \emph{lakṣyam}
  \item \emph{Say it:} \emph{lakṣyam}, once more
  \item \emph{Say it:} say \emph{lakṣyam}
  \item \emph{Recall:} say the Sanskrit for a dispute, then the Sanskrit for a plan, then say \emph{lakṣyam} again
\end{itemize}

\begin{tcolorbox}[breakable,colback=teal!4,colframe=teal!35!black,title={Before you move on}]
What does \sk{लक्ष्यम्} mean? (\textquotedblleft{}A goal, a target\textquotedblright{}.) Say it once more.
\end{tcolorbox}
\section[\texorpdfstring{\sk{उपायः} (upāyaḥ)}{उपायः (upāyaḥ)}]{\sk{उपायः} (upāyaḥ) — a way, a means}

\label{lesson:SA-C201-upayah}



\begin{quote}
Before the new one: say the Sanskrit for a plan, then the Sanskrit for a goal.
\end{quote}

\subsection*{You'll want to know: \sk{उपायः}}
\textbf{\sk{उपायः}} — \emph{upāyaḥ} — \textquotedblleft{}a way, a means\textquotedblright{}.

\begin{grammarlens}[title={What you've built}]
One more word for this chapter.
\end{grammarlens}

\subsection*{Your turn}
\begin{itemize}
\raggedright
  \item \emph{Say it:} \emph{upāyaḥ}
  \item \emph{Say it:} \emph{upāyaḥ}, once more
  \item \emph{Say it:} say \emph{upāyaḥ}
  \item \emph{Recall:} say the Sanskrit for a plan, then the Sanskrit for a goal, then say \emph{upāyaḥ} again
\end{itemize}

\begin{tcolorbox}[breakable,colback=teal!4,colframe=teal!35!black,title={Before you move on}]
What does \sk{उपायः} mean? (\textquotedblleft{}A way, a means\textquotedblright{}.) Say it once more.
\end{tcolorbox}
\section[\texorpdfstring{\sk{समस्या} (samasyā)}{समस्या (samasyā)}]{\sk{समस्या} (samasyā) — a problem}

\label{lesson:SA-C201-samasya}



\begin{quote}
Before the new one: say the Sanskrit for a goal, then the Sanskrit for a way.
\end{quote}

\subsection*{You'll want to know: \sk{समस्या}}
\textbf{\sk{समस्या}} — \emph{samasyā} — \textquotedblleft{}a problem\textquotedblright{}.

\begin{grammarlens}[title={What you've built}]
One more word for this chapter.
\end{grammarlens}

\subsection*{Your turn}
\begin{itemize}
\raggedright
  \item \emph{Say it:} \emph{samasyā}
  \item \emph{Say it:} \emph{samasyā}, once more
  \item \emph{Say it:} say \emph{samasyā}
  \item \emph{Recall:} say the Sanskrit for a goal, then the Sanskrit for a way, then say \emph{samasyā} again
\end{itemize}

\begin{tcolorbox}[breakable,colback=teal!4,colframe=teal!35!black,title={Before you move on}]
What does \sk{समस्या} mean? (\textquotedblleft{}A problem\textquotedblright{}.) Say it once more.
\end{tcolorbox}
\section[\texorpdfstring{\sk{समाधानम्} (samādhānam)}{समाधानम् (samādhānam)}]{\sk{समाधानम्} (samādhānam) — a solution}

\label{lesson:SA-C201-samadhanam}



\begin{quote}
Before the new one: say the Sanskrit for a way, then the Sanskrit for a problem.
\end{quote}

\subsection*{You'll want to know: \sk{समाधानम्}}
\textbf{\sk{समाधानम्}} — \emph{samādhānam} — \textquotedblleft{}a solution\textquotedblright{}.

\begin{grammarlens}[title={What you've built}]
One more word for this chapter.
\end{grammarlens}

\subsection*{Your turn}
\begin{itemize}
\raggedright
  \item \emph{Say it:} \emph{samādhānam}
  \item \emph{Say it:} \emph{samādhānam}, once more
  \item \emph{Say it:} say \emph{samādhānam}
  \item \emph{Recall:} say the Sanskrit for a way, then the Sanskrit for a problem, then say \emph{samādhānam} again
\end{itemize}

\begin{tcolorbox}[breakable,colback=teal!4,colframe=teal!35!black,title={Before you move on}]
What does \sk{समाधानम्} mean? (\textquotedblleft{}A solution\textquotedblright{}.) Say it once more.
\end{tcolorbox}
